\chapter{Determining the Class Group $C_{K}$}
Recall $C_{K}=\{[I]|I\lhd \mathcal{O}_{K} \text{ nonzero ideal}\}$ is a finite abelian group. $\left| C_{K} \right| =1$ if and only if $\mathcal{O}_{K}$ is a PID. We will study the structure of $C_{K}$ and use it to solve Diophantine equations. 

\section{Minkowski's Convex Body Theorem}
Let $\mathscr{B}=\{\mathbf{v}_{1},\dots,\mathbf{v}_{n}\}$ be an $\mathbb{R}$-basis of $\mathbb{R}^{n}$ and let 
\[
L:=\{\sum_{i=1}^{n}a_{i}\mathbf{v}_{i}|a_{i}\in \mathbb{Z}\}
\]
be a {\bf lattice} generated by $\mathscr{B}$. 
The {\bf fundamental domain} of $L$ is 
\[
D=\{\sum_{i=1}^{n}\alpha_{i}\mathbf{v}_{i}|0\leq \alpha_{i}<1\}. 
\]
\begin{figure}[h]
    \centering
    \begin{subfigure}[b]{0.45\textwidth}
        \centering
        \includegraphics[width=\textwidth]{figures/Figure_1.png}
        \caption{Fundamental Domain \(D\)}
        \label{fig:first}
    \end{subfigure}
    \begin{subfigure}[b]{0.45\textwidth}
        \centering
        \includegraphics[width=\textwidth]{figures/Figure_2.png}
        \caption{Example of \(S(\mathbf{a})\): Yellow Area}
        \label{fig:second}
    \end{subfigure}
    \caption{\(\mathbb{R}^2\) Example}
    \label{fig:parallel_images}
\end{figure}

For all $\mathbf{v}\in \mathbb{R}^{n}$, $\mathbf{v}$ is uniquely expressed as $\mathbf{v}=\mathbf{u}+\mathbf{w}$ where $\mathbf{u}\in L$ and $\mathbf{w}\in D$. Define the volume of $L$ by $vol(L):=vol(D)=\left| \det(\mathbf{v}_{1},\dots,\mathbf{v}_{n})\right|$. 



\begin{proposition}
    Let $L$ be a lattice in $\mathbb{R}^{n}$, $S\subseteq \mathbb{R}^{n}$ be a bounded and measurable set such that $vol(S)>vol(L)$. Then $\exists \mathbf{x}, \mathbf{y}\in S$ such that $\mathbf{x}\neq \mathbf{y}$ and $\mathbf{x}-\mathbf{y}\in L$. 
\end{proposition}

\begin{proof}
    For all $\mathbf{a}\in L$, let $S(\mathbf{a})=(S-\mathbf{a})\cap D$. Then $S=\coprod\limits_{\mathbf{a}\in L}S(\mathbf{a})+\mathbf{a}$. Since 
    \[
    vol(S)=\sum_{\mathbf{a}\in L}vol(S(\mathbf{a}))>vol(L),
    \]
    there must be some $\mathbf{a}\neq \mathbf{b}\in L$ such that $S(\mathbf{a})\cap S(\mathbf{b})\neq \varnothing$. We may find $\mathbf{x}=\mathbf{v}+\mathbf{a}$, $\mathbf{y}=\mathbf{v}+\mathbf{b}\in S$, so $\mathbf{x}-\mathbf{y}=\mathbf{a}-\mathbf{b}\in L$. 
\end{proof}

\begin{theorem}[Minkowski's Convex Body]\label{thm-convexBody}
    Let $L$ be a lattice in $\mathbb{R}^{n}$. $S\subseteq \mathbb{R}^{n}$ be a bounded, measurable, convex ($\forall \mathbf{x},\mathbf{y}\in S$ and $\forall 0\leq t\leq 1$, the linear combination $t\mathbf{x}+(1-t)\mathbf{y}\in S$) and symmetric ($\mathbf{x}\in S\Rightarrow -\mathbf{x}\in S$). Suppose $vol(S)>2^{n}vol(L)$, then there exists a nonzero $\mathbf{v}\in S\cap L$. 
\end{theorem}

\begin{proof}
    Consider $vol(\frac{1}{2}S)=\frac{1}{2^{n}}vol(S)$. Then by above proposition, there exists $\mathbf{x},\mathbf{y}\in \frac{1}{2}S$ such that $\mathbf{x}-\mathbf{y}\in L$. Since $S$ is symmetric, $2\mathbf{x},-2\mathbf{y}\in S$. By the convexity of $S$, $\mathbf{x}-\mathbf{y}=\frac{1}{2}(2\mathbf{x})+\frac{1}{2}(-2\mathbf{y})\in S$. Since $\mathbf{x}\neq\mathbf{y}$, this is a nonzero element of $S\cap L$. 
\end{proof}

\begin{remark}
    If $S$ is closed, then one can replace $vol(S)>2^{n}vol(L)$ by $vol(S)\geq 2^{n}vol(L)$. To see so, apply the theorem to $(1+\frac{1}{n})S$. Then there exist $\mathbf{v}_{n}\in (1+\frac{1}{n})S\cap L$. Since 
    \[
    \bigcap_{n=1}^{\infty}(1+\frac{1}{n})S=S
    \] 
    and 
    \[
    S\subseteq \cdots\subseteq (1+\frac{1}{n})S\subseteq 2S, 
    \]
    the sequence $(\mathbf{v}_{n})$ has a subsequence convergent to some $\mathbf{v}\in S$. 
\end{remark}

\begin{example}
    Let $p=1 \mod 4$ be a rational prime, then $\exists x,y\in \mathbb{Z}$ s.t. $x^{2}+y^{2}=p$. 
    \begin{proof}
        Let $s\in \mathbb{Z}_{p}^{*}$ be an element of order 4, such element exists since $\mathbb{Z}_{p}^{*}$ is an cyclic group of order $p-1$. Then $s^{2}=-1 \mod p$. 

        For $x,y \in \mathbb{Z}$ such that $x^{2}+y^{2}=0 \mod p$, then $(x+sy)(x-sy)=0\mod p$, i.e. $x=\pm sy\mod p$. By switching $s\to -s$ if necessary, we may assume $x=sy\mod p$. So $x=sy+zp$ for some $z\in \mathbb{Z}$. We may write 
        \[
        (x,y)=(sy+zp,y)=(s,1)y+(p,0)z. 
        \]
        We only need to find such $(x,y)$. 
        Consider the lattice $L=\langle (s,1),(p,0)\rangle_{\mathbb{Z}}$, then $vol(L)=p$. Let $S=\{(x,y)|x^{2}+y^{2}<2p\}$, then $vol(S)=2\pi p>2^{2}vol(L)$. By Theorem \ref{thm-convexBody}, there exists $(x,y)\in S\cap L$, i.e. $x^{2}+y^{2}<2p$ and $x=sy\mod p$. The only possible case is $x^{2}+y^{2}=p$. 
    \end{proof}
\end{example}

\section{Minkowski's Bound} \label{sec-minkowski}
In this section, we will find a lattice and a bounded, measurable, convex and symmetric set and we will use Theorem \ref{thm-convexBody} to prove Theorem \ref{thm-MinkowskiBound}. 

Let $K:\mathbb{Q}$ be a number field with $[K:\mathbb{Q}]=n=r+2s$, where 
\[
\sigma_{1},\dots, \sigma_{r}:K\hookrightarrow\mathbb{R} \quad\text{real embeddings}
\]
\[
\left.\begin{aligned}
    \sigma_{r+1},\dots, \sigma_{r+s}\\
    \overline{\sigma_{r+1}},\dots, \overline{\sigma_{r+s}}
\end{aligned}\right\}:K\hookrightarrow \mathbb{C}\quad \text{complex embeddings}
\]

\begin{example}
    \begin{enumerate}
        \item Let $K=\mathbb{Q}(\sqrt{-5})$, then 
        \[
        \left.\begin{aligned}
            \sigma_{1}:a+b\sqrt{-5}\mapsto a+b\sqrt{-5}\\
            \sigma_{2}:a+b\sqrt{-5}\mapsto a-b\sqrt{-5}
        \end{aligned}\right\} \quad \text{are all complex embeddings}
        \]
        \item Let $K=\mathbb{Q}(\sqrt{3})$, then  
        \[
        \left.\begin{aligned}
            \sigma_{1}:a+b\sqrt{3}\mapsto a+b\sqrt{3}\\
            \sigma_{2}:a+b\sqrt{3}\mapsto a-b\sqrt{3}
        \end{aligned}\right\} \quad \text{are all real embeddings}
        \]
        \item Let $K=\mathbb{Q}(\sqrt[3]{2})$, then 
        \[
        \sigma_{1}:\sqrt[3]{2}\mapsto \sqrt[3]{2}\quad \text{ is real embedding}
        \]
        \[
        \left.\begin{aligned}
            \sigma_{2}:\sqrt[3]{2}\mapsto \sqrt[3]{2}\omega\\
            \overline{\sigma_{2}}:\sqrt[3]{2}\mapsto \sqrt[3]{2}\omega^{2}
        \end{aligned}\right\} \quad \text{are complex embeddings}
        \]
    \end{enumerate}
\end{example}

Define $\sigma:K\to \mathbb{R}^{r}\times \mathbb{C}^{s}\cong \mathbb{R}^{n}$ given by 
\[
\sigma: k\mapsto \begin{pmatrix}
    \sigma_{1}(k)\\
    \vdots\\
    \sigma_{r}(k)\\
    \mathrm{Re}(\sigma_{r+1}(k))\\
    \mathrm{Im}(\sigma_{r+1}(k))\\
    \vdots\\
    \mathrm{Re}(\sigma_{r+s}(k))\\
    \mathrm{Im}(\sigma_{r+s}(k))\\
\end{pmatrix}
\]

\begin{remark}
    Let $\{k_{1},\dots, k_{n}\}\subseteq K$ be a $\mathbb{Q}$-basis of $K$ (e.g. integral basis $\{\omega_{1},\dots,\omega_{n}\}$ of $\mathcal{O}_{K}$). Then 
    \[
    \begin{aligned}
        \det\begin{pmatrix} \uparrow & & \uparrow\\ \sigma(k_{1}) & \dots & \sigma(k_{n}) \\
        \downarrow & & \downarrow \end{pmatrix}&=\det\begin{pmatrix}
            \sigma_{1}(k_{1}) & &&  &\sigma_{1}(k_{n}) \\
            \vdots &&&&\vdots \\
            \sigma_{r}(k_{1})&&&&\sigma_{r}(k_{n})\\
            \mathrm{Re}(\sigma_{r+1}(k_{1}))&&&&\mathrm{Re}(\sigma_{r+1}(k_{n}))\\
            \mathrm{Im}(\sigma_{r+1}(k_{1}))&&\cdots&&\mathrm{Im}(\sigma_{r+1}(k_{n}))\\
            \vdots&&&&\vdots\\
            \mathrm{Re}(\sigma_{r+s}(k_{1}))&&&&\mathrm{Re}(\sigma_{r+s}(k_{n}))\\
            \mathrm{Im}(\sigma_{r+s}(k_{1}))&&&&\mathrm{Im}(\sigma_{r+s}(k_{n}))\\
        \end{pmatrix}\\
        &=\frac{1}{(2i)^{s}}\det \begin{pmatrix}
            \sigma_{1}(k_{1})& &&&\sigma_{1}(k_{n})\\
            \vdots &&&&\vdots \\
            \sigma_{r}(k_{1})&&& &\sigma_{r}(k_{n})\\
            \sigma_{r+1}(k_{1})& &\cdots& &\sigma_{r+1}(k_{n})\\
            \overline{\sigma_{r+1}}(k_{1})&&&&\overline{\sigma_{r+1}}(k_{n})\\
            \vdots&&&&\vdots\\
            \overline{\sigma_{r+s}}(k_{1})&&&&\overline{\sigma_{r+s}}(k_{n})
        \end{pmatrix}
    \end{aligned}
    \]
    (the last equality comes from simple elementary row operation $\begin{pmatrix}
        x \\ y
    \end{pmatrix} \to \begin{pmatrix}
        x+yi \\ x-yi
    \end{pmatrix}$). 
    In particular, 
    \[\det\begin{pmatrix} \uparrow & & \uparrow\\ \sigma(\omega_{1}) & \dots & \sigma(\omega_{n}) \\
        \downarrow & & \downarrow \end{pmatrix}=\frac{1}{(2i)^{s}}\det(\sigma_{i}(\omega_{j}))=\frac{1}{(2i)^{s}}\Delta(\underline{\omega})=\frac{1}{(2i)^{s}}\sqrt{\Delta^{2}(K)}.\]
\end{remark}

Now we consider $I\lhd \mathcal{O}_{K}$ a nonzero ideal. Note that $[\mathcal{O}_{K}:I]=N(I)<\infty$. $\mathcal{O}_{K}=\langle \omega_{1},\dots, \omega_{n}\rangle_{\mathbb{Z}}$ where $\omega_{1},\dots,\omega_{n}$ is the integral basis (also a $\mathbb{Q}$-basis of $K$). Then $\mathcal{O}_{K}$ forms a lattice generated by $\{\sigma(\omega_{1}),\dots,\sigma(\omega_{n})\}$. 

$I$ is a sublattice of finite index so $I=\langle k_{1},\dots,k_{n}\rangle_{\mathbb{Z}}$, where $\{k_{1},\dots,k_{n}\}$ is also a $\mathbb{Q}$-basis of $K$. Denote the sublattice correspond to $I$ by $L_{I}$, which is generated by $\{\sigma(k_{1}),\dots, \sigma(k_{n})\}$. By the previous remark, 
\[
\begin{aligned}
    vol(L_{I})&=\left| \det(\sigma(k_{1}),\dots, \sigma(k_{n})) \right|=\frac{1}{2^{s}}\sqrt{\left| \Delta^{2}(k_{1},\dots, k_{n}) \right| }\\
    &=\frac{1}{2^{s}}\sqrt{[\mathcal{O}_{K}:I]^{2}\left| \Delta^{2}(K) \right| }=\frac{N(I)}{2^{s}}\sqrt{\left| \Delta^{2}(K) \right| }. 
\end{aligned}
\]
This $L_{I}$ will serve as the lattice.

For $t>0$, define set 
\[R_{t}=\{\begin{pmatrix}
    x_{1}\\
    \vdots\\
    x_{r}\\
    z_{r+1}\\
    \vdots\\
    z_{r+s}
\end{pmatrix}\in\mathbb{R}^{r}\times \mathbb{C}^{s}\cong \mathbb{R}^{n}|\sum_{i=1}^{r}\left| x_{i} \right| +\sum_{j=1}^{s}2\left| z_{r+j} \right| \leq t\}. \]
Then clearly $R_{t}$ is compact, convex and symmetric, whose volume is
\(
vol(R_{t})=\frac{2^{r}t^{n}(\frac{\pi}{2})^{s}}{n!}.
\)

\begin{theorem}\label{thm-MinkowskiBound}
    Let $I\lhd \mathcal{O}_{K}$ be a nonzero ideal, then $\exists \alpha\in I\backslash \{0\}$ such that 
    \[
    \left| Norm_{K /\mathbb{Q}} (\alpha)\right| \leq m_{K}N(I),
    \]
    where $m_{K}=(\frac{4}{\pi})^{s}\frac{n!}{n^{n}}\sqrt{\left| \Delta^{2}(K) \right| }$. 
\end{theorem}

\begin{proof}
    Take $t>0$ such that $\frac{\pi^{s}t^{n}}{n!}=4^{s}\sqrt{\left| \Delta^{2}(K) \right| }N(I)$ and consider $R_{t}\subseteq \mathbb{R}^{n}$. 
    \[vol(R_{t})=\frac{2^{r}t^{n}(\frac{\pi}{2})^{s}}{n!}=2^{r+2s}vol(L_{I})=2^{n}vol(L_{I}).\]
    By Theorem \ref{thm-convexBody}, there exists $\alpha\in I\backslash \{0\}$ s.t. $\sigma(\alpha)\in R_{t}$, i.e. 
    \[
    \sum_{i=1}^{r}\left| \sigma_{i}(\alpha) \right| +\sum_{j=1}^{s}2\left| \sigma_{r+j}(\alpha) \right|\leq t. 
    \]
    So 
    \[
    \frac{1}{n}\left(\sum_{i=1}^{r}\left| \sigma_{i}(\alpha) \right| +\sum_{j=1}^{s}(\left| \sigma_{r+j}(\alpha) \right|+\left| \overline{\sigma_{r+j}}(\alpha)\right|)\right) \leq \frac{t}{n},
    \]
    by AM-GM inequality, 
    \[\sqrt[n]{\prod_{\sigma\in \mathrm{Hom}(K,\mathbb{C})}\left| \sigma(\alpha) \right| }\leq \frac{t}{n}. \]
    Hence 
    \[
    \left| Norm_{K /\mathbb{Q}}(\alpha)\right|\leq (\frac{t}{n})^{n}=(\frac{4}{\pi})^{s}\frac{n!}{n^{n}}\sqrt{\left| \Delta^{2}(K) \right| }N(I). 
    \]
\end{proof}

\begin{corollary}
    Let $[K:\mathbb{Q}]=n$ be a number field. For any $c\in C_{K}=\{[I]|I\lhd\mathcal{O}_{K}\text{ nonzero}\}$, there exists $J\lhd \mathcal{O}_{K}$ such that $c=[J]$ and $N(J)\leq m_{K}$. 
\end{corollary}

\begin{proof}
    Fix $c\in C_{K}$ and consider $I\lhd \mathcal{O}_{K}$ such that $[I]=c^{-1}$. By Theorem \ref{thm-MinkowskiBound}, there exists $\alpha\in I\backslash \{0\}$ such that $\left| Norm_{K /\mathbb{Q}}(\alpha) \right| \leq m_{K}N(I)$. Note that 
    \begin{center}
    $(\alpha)\subseteq I$\quad  $\Rightarrow$ \quad $I\mid (\alpha)$\quad  $\Rightarrow$ \quad  there exists $J\lhd \mathcal{O}_{K}$ such that $IJ=(\alpha)$.    
    \end{center}
    So 
    \[
    N(I)N(J)=N(IJ)=N((\alpha)) = \left|Norm_{K /\mathbb{Q}}(\alpha)\right|\leq m_{K}N(I)
    \] 
    which gives $N(J)\leq m_{K}$. Moreover, since $IJ = (\alpha)$, $[I]*[J] = [(\alpha)] = e$ and hence $[J]=[I]^{-1}=c$. So the result follows.
\end{proof}

As an immediate consequence of the above corollary, one has $C_{K}=\{[J]\ |\ N(J) \leq m_{K}\}$. In other words, we can list all elements in $C_{K}$ by finding all ideals $J$ of norm less than or equal to $m_K$.

\section{Calculating Class Groups}
\begin{example}
    Let $K=\mathbb{Q}(\sqrt{-5})$, $\mathcal{O}_{K}=\mathbb{Z}[\sqrt{-5}]$. Then $r=0$ and $s=1$, $\left| \Delta^{2}(K) \right| =20$, so $m_{K}=\frac{4\sqrt{5}}{\pi} \approx 2.\cdots$. So 
    $$C_{K}=\{[J]\ |\ N(J)\leq 2\}.$$ 

    So we only have to study ideals with $N(J) = 1$ or $2$. In the case $N(J)=1$, $J=\mathcal{O}_{K}$ and hence $[J] = e$ is the identity element in $C_K$. As for $N(J)=2$, $[\mathcal{O}_{K}:J]=2$, so $J\lhd \mathcal{O}_{K}$ must be a prime ideal lying above 2. By Theorem \ref{thm-factorizationOfRationalPrimes}, the only possibility is the non-principal ideal $J=(2,1+\sqrt{-5})$. 
    Therefore, 
    $$C_K = \{e:=[\mathcal{O}_K],\ \  \tau:=[(2,1+\sqrt{-5})]\}$$ 
    has only two elements, i.e. $C_{K}\cong \mathbb{Z}_{2}$. Alternatively, one can check directly that
    $\tau^{2}=[(2,1+ \sqrt{-5})^2] = [(2)] =e$. 
\end{example}

\begin{example}
    Let \(K = \mathbb{Q}(\sqrt{-6}), \mathcal{O}_K = \mathbb{Z}[\sqrt{-6}] \). One can compute \(m_K = \frac{2}{\pi}\sqrt{24} \approx 3.\cdots\). Therefore
    \begin{align*}
        C_K = \left\{[J]\ |\  N(J)= 1,2,3\right\}
    \end{align*}
    Note that
    \begin{align*}
        N(J_1) = 1 &\Leftrightarrow J_1 = \mathcal{O}_K\\
        N(J_2) = 2 &\Leftrightarrow J_2 \triangleleft \mathcal{O}_K \text{ is a maximal ideal above }2. \text{ Since } (2) = (2,\sqrt{-6})^2, \text{ this implies } J_2 = (2,\sqrt{-6})\\
        N(J_3) = 3 &\Leftrightarrow J_3 = (3,\sqrt{-6}).
    \end{align*}
    Therefore, \[C_K = \left\{e=[\mathcal{O}_K], [J_2], [J_3]\right\}.\] 
    We can easily check that \([J_2] \neq e\) is not principal, otherwise suppose \(J_2 = (a+b\sqrt{-6})\) for some \(a,b \in \mathbb{Z}\), \(N(J_2)= a^2+6b^2 = 2\). The only choice for \(b\) is \(b = 0\), implies  \(a^2 =2\), \(a \notin \mathbb{Z}\) contradiction. Similarly one has \([J_3] \neq e\). Also, note that 
    \[[J_2]^2 = [(2)], \quad [J_3]^2 = [(3)]  \quad \quad \Rightarrow \quad \quad [J_2]^2 = [J_3]^2 = e.\]
    which implies \(C_K\) has no elements of order 3. Thus \(C_K \ncong \mathbb{Z}_3\), implying that $[J_2] = [J_3]$ and 
    \[C_K = \{e, [J_2] = [J_3]\} \cong \mathbb{Z}_2.\]

    Alternatively, one can also check directly that $[J_2] = [J_3]$: since \(\sqrt{-6} = 3\sqrt{-6} + \sqrt{-6}(-2)\) we have
    \begin{align*}
        \sqrt{-6} \in J_3 J_2 \quad \Rightarrow \quad (\sqrt{-6}) \subseteq J_3 J_2\quad \Rightarrow \quad J_3 J_2 | (\sqrt{-6})
    \end{align*}
    By computing norm \(N(J_2) N(J_3) = 2 \times 3 = 6 = N\left((\sqrt{-6})\right)\), we can conclude that 
    \begin{align*}
        J_3J_2 = (\sqrt{-6}) \quad \Rightarrow \quad [J_3]*[J_2] = [(\sqrt{6})] = e \quad \Rightarrow \quad [J_3] = [J_2]^{-1} = [J_2]
    \end{align*}
    since $[J_2]^2 = e$.
\end{example}

\begin{remark} \label{rmk-ckgenerator}
    By Minikowski Bound, we can easily have \begin{align*}
        C_K = \left\{[J]| N(J) \leq m_K\right\}.
    \end{align*}
    If we only care about the generators of \(C_K\) (as a finite abelian group), we only need to study \(J\) with \(N(J) = P\) for \(p \leq m_K\) rational prime, i.e.
    $$C_K = \langle J\ |\ N(J) = p,\ p \leq m_K\ \text{rational prime}\rangle.$$
    For instance, suppose \(N(J) = 6\). Then for all \(\alpha \in \mathcal{O}_K, 6\alpha \in J\) (since \(N(J) = |\mathcal{O}_K/J|=6\)). We have
    \begin{align*}
        (6) \subseteq J \quad \Rightarrow \quad J|(6)\quad \Rightarrow \quad J|(2)(3).
    \end{align*}
    By Dedekind's theorem, \(J\) can be factorized into (some prime ideals lying above 2) \(\times\) (some prime ideals lying above 3). In other words, the prime ideals lying above $2$ and $3$ generates \(J\).
\end{remark}

\begin{example}
    Let \(K = \mathbb{Q}(\sqrt{-163})\) and \(\mathcal{O}_K = \displaystyle \mathbb{Z}[\frac{1+\sqrt{-163}}{2}]\). One can easily compute \(m_K = \frac{4}{\pi} \frac{1}{2} \sqrt{163} = 8. \cdots\).

    By the Remark above, 
    \[C_K = \langle J\ |\ N(J) = 2,3,5,7 \rangle\] 
    So we can apply theorem \ref{thm-factorizationOfRationalPrimes}, with \(\alpha = \frac{1+\sqrt{-163}}{2}\) to find the ideals $J$ lying above $2,3,5,7$. Since the minimal polynomial of $\alpha$ is
    \begin{align*}
        m_\alpha(x) = x^2 - x + 41,
    \end{align*}
    we can check that 
    \begin{itemize}
        \item[\((2)\):] \(\bar m_\alpha(x) = x^2 + x + 1\) irreducible in \(\mathbb{Z}_2[x] \Rightarrow (2)\) is prime;
        \item[\((3)\):] \(\bar m_\alpha(x) = x^2 + x - 1\) irreducible in \(\mathbb{Z}_3[x] \Rightarrow (3)\) is prime;
        \item[\((5)\):] \(\bar m_\alpha(x) = x^2 + x + 1\) irreducible in \(\mathbb{Z}_5[x] \Rightarrow (5)\) is prime;
        \item[\((7)\):] \(\bar m_\alpha(x) = x^2 + x - 1\) irreducible in \(\mathbb{Z}_7[x] \Rightarrow (7)\) is prime.
    \end{itemize}
    Therefore there is no \(J\) with \(N(J) =2,3,5,7\). Hence 
    \begin{align*}
        C_K = <\emptyset> = {e}
    \end{align*}
    is the trivial group. \(i.e.\) \(\mathcal{O}_K\) is a PID.
\end{example}

\begin{example}
    Let \(K = \mathbb{Q}(\sqrt{-29})\) and \(\mathcal{O}_K = \displaystyle \mathbb{Z}[-29]\). One can easily compute \(m_K = \frac{2}{\pi} \sqrt{116} = 6. \cdots\). So we need to study \(J \triangleleft \mathcal{O}_K\) with \(N(J) = 2,3,5\), where
    \begin{align*}
        (2) &=  P^2 = (2,1+\sqrt{-29})^2;\\
        (3) &= Q Q' = (3, 1+\sqrt{-29})(3,1-\sqrt{-29});\\
        (5) &= R R' = (5, 1+\sqrt{-29})(5, 1-\sqrt{-29}).
    \end{align*}
    we can get \(C_K = \left< [P], [Q], [Q'], [R], [R']  \right>\). We want to study the relation between ideal classes \([P],[Q],[R], [Q'],[R']\). Obviously we have 
    \begin{align*}
        [P]^2 = [Q][Q'] = [R][R'] = e\\
        [Q'] = [Q]^{-1} = [Q]^a\\
        [R'] = [R]^{-1} = [R]^b
    \end{align*}
    for some \(a,b > 1\) (otherwise $Q' = Q$ or $R' = R$, contradicting the unique factorization in Theorem \ref{thm-factorizationOfRationalPrimes}). We can simplify \(C_K = \left< [P], [Q], [R]\right>\), where \([P]\) have order 2 and \([Q]\) and $[R]$ has order greater or equal to \(3\).
    
    To look at the orders of $[Q]$ and $[R]$, one needs to do some trial and error. By observation, $\mathrm{Norm}_{K/\mathbb{Q}}(3+2\sqrt{-29}) = 3^2 + 4\cdot 29 = 125 = 5^3$, so we study the prinicpal ideal
    \(I = (3+2 \sqrt{-29})\). 

    By the argument given in Remark \ref{rmk-ckgenerator}, 
    $$I\ |\ (125) = (5)^3 = R^3(R')^3,$$
    thus there are 4 possible ways in prime factorization as follows
    \begin{align*}
        I = R^3, R^2 R' , R (R')^2 , (R')^3.
    \end{align*}
    Suppose \(R|I\), implying that \(I \subseteq R\), we have
    \begin{align*}
        3+2\sqrt{-29} \in I \subseteq R,\quad \quad 2 (1+\sqrt{-29}) \in R,\\ 
    \end{align*}
    Therefore, 
    \begin{align*}
        \left\{(3+2\sqrt{-29}) - 2(1+\sqrt{-29}) \right\} \in R\quad 
        \Rightarrow \quad 1\in R\quad
        \Rightarrow \quad R= \mathcal{O}_K,
    \end{align*}
    which is a contradiction. Therefore the only choice for prime factorization is \(I = (R')^3\), and we can compute
\begin{align*}
    [R']^3 = [I] =e \quad 
    \Rightarrow\quad  \left([R]^{-1}\right)^3 = e\quad
    \Rightarrow \quad [R]^3 = e.
\end{align*}
We conclude that degree of \([R]\) is \(3\). 

\medskip
Now we study \(I' = (1+\sqrt{-29})\) with \(N(J) = 30\), one can easily check that \(I' | (2)(3)(5)\), which implies that 
\begin{align*}
    I' | P^2\cdot Q Q'  \cdot RR'.
\end{align*}
Since \(N(I') = 30 = 2\times3 \times 5\), therefore we have
\begin{align*}
    I = \left\{\begin{matrix}
PQR \\
PQ'R \\
PQR' \\
PQ'R'
\end{matrix}\right. \quad \Rightarrow \quad [I] = \left\{\begin{matrix}
[P][Q][R] \\
[P][Q]^{-1}[R] \\
[P][Q][R]^{-1} \\
[P][Q]^{-1}[R]^{-1}
\end{matrix}\right.
\end{align*}
hence \([Q]\) can be expressed as
\begin{align*}
    [Q] =  [P][R] \quad \text{or}\quad [P][R]^{-1} \in \left<[P] ,[R]\right>.
\end{align*}

As a conclusion, 
\[C_K = \langle [P],[R]\ |\ [P]^2 = e,\ [R]^3 = e \rangle \cong \mathbb{Z}_2 \times \mathbb{Z}_3 \cong \mathbb{Z}_6.\]
\end{example}


\section{Applications to Diophantine Equations}
We will end by presenting two examples on how the knowledge on $C_K$ helps solve certain Diophantine equations:
\begin{example}
Show that there are no integer solutions to $Y^2+37 = X^3$.
\end{example}
\begin{proof}
Let $K = \mathbb{Q}(\sqrt{37})$, with $\mathcal{O}_K = \mathbb{Z}[\sqrt{-37}]$. One can check that $C_K \cong \mathbb{Z}_2$. Then the equation can be factorized as:
$$\{Y + \sqrt{-37}\}\{Y - \sqrt{-37}\} = X^3$$
\noindent {\bf Claim 1. $2, 37 \nmid X$}. 

To see so, suppose on contrary that $2 \mid X$, then $X^3 = Y^2 + 37 \equiv 0 \mod 8$, which implies $Y^2 \equiv 3 \mod 8$. However, one can easily check that all square numbers can only have $n^2 \equiv 0, 1, 4 \mod8 $.

On the other hand, if $37 \mid X$, then obviously one also has $37 \mid Y$. Write $X = 37a$ and $Y = 37b$, then
$$(37b)^2 + 37 = (37a)^3 \quad \Rightarrow \quad 37b^2 + 1 = 37^2a^3 \equiv 0 \mod 37$$
But $37b^2 + 1 \equiv 1 \mod 37$, which again gives a contradiction. \qed

\medskip 
\noindent{\bf Claim 2. The principal ideals $(Y + \sqrt{-37})$ \& $(Y - \sqrt{-37})$ are coprime}. 

Suppose on contrary there exists a prime ideal $P \mid (Y + \sqrt{-37})$ and $ P \mid (Y - \sqrt{-37})$. Thus 
$$(Y \pm \sqrt{-37}) \subseteq P \quad \Rightarrow  \quad Y \pm \sqrt{-37} \in P  \quad \Rightarrow  \quad 2\sqrt{-37} = \{Y + \sqrt{-37}\} - \{Y - \sqrt{-37}\} \in P$$  
and hence
$$(2\sqrt{-37}) \subseteq P  \quad
\Rightarrow  \quad P \mid (2\sqrt{-37}) = (2)(\sqrt{-37})  \quad \Rightarrow  \quad P \mid (2) \text{ or } P \mid (\sqrt{-37})$$
Therefore, $P = (2, 1+ \sqrt{-37})$ or $(\sqrt{-37})$. 

Suppose $P = (2, 1+ \sqrt{-37})$, then 
\begin{align*}
(2) = P^2 | (Y+\sqrt{-37})(Y-\sqrt{-37}) = (Y^2+37) = (X^3) \quad &\Rightarrow \quad N((2)) \mid N((X^3)) \\
&\Rightarrow \quad  4 \mid X^6 \\ &\Rightarrow \quad 2 \mid X,\end{align*}
contradicting Claim 1.

On the other hand, suppose $P = (\sqrt{-37})$, then
\begin{align*}
Y + \sqrt{-37} \in (\sqrt{-37}) \ &\Rightarrow \ Y = \{Y + \sqrt{-37}\} - \sqrt{-37} \in (\sqrt{-37}) \\ &\Rightarrow \ (\sqrt{37}) \mid (Y) \\
&\Rightarrow \  37 = N((\sqrt{37})) \mid N((Y)) = Y^2 
\end{align*}
Hence $37 \mid Y$, and the original equation implies that $37 \mid X$, contradicting Claim 1 again.  \qed

\bigskip
Consider the ideal factorization
$$(Y + \sqrt{-37})(Y - \sqrt{-37}) = (X)^3 = P_1^{3e_1} P_2^{3e_2} \cdots P_r^{3e_r}$$ 
where $(X) = P_1^{e_1} P_2^{e_2}\dots P_r^{e_r}$ is the prime ideal factorization of $(X)$. By Claim 2 and Dedekind's Theorem, if the prime ideal $P_i \mid (Y+\sqrt{-37})$, then
$P_i \nmid (Y-\sqrt{-37})$ and hence $P_i^{3e_i} \nmid (Y+\sqrt{-37})$. Therefore,
$$(Y + \sqrt{-37}) = I^3  \quad \quad (*)$$ 
where $I = P_{i_1}^{e_{i_1}}\dots P_{i_l}^{e_{i_l}}$.

Now study ideal class group:
$$e = [(Y + \sqrt{-37})] = [I]^3$$
Since $|C_K| = 2$, $[I]^2 = e$ and hence $[I]^3 = [I]$, which gives
$$e = [(Y + \sqrt{-37})] = [I],$$
i.e. $I = (a+b\sqrt{-37})$ itself is a principal ideal. Hence $(*)$ gives
\begin{align*}
(Y + \sqrt{-37}) &= (a+b\sqrt{-37})^3 = (\{a^3-111ab^2\} + \{3a^2b-37\}\sqrt{-37}) \\ 
\Rightarrow \quad Y + \sqrt{-37} &= \pm\{\{a^3-111ab^2\} + \{3a^2b-37b^3\}\sqrt{-37}\}
\end{align*}

Since $\{1,\sqrt{-37}\}$ is a $\mathbb{Q}$-basis of $K$, they are linearly independent and hence
$$Y = \pm a(a^2-111b^2), \quad \quad  1 = \pm b(3a^2-37b^2).$$
But the second equation implies that $b = \pm 1$ and hence
$3a^2 - 37 = \pm 1$, which implies 
$3a^2 = 38$ or $36$. But this is impossible.

\medskip
Consequently, the original equation $Y^2+37 = X^3$ has no integer solutions.
\end{proof}



\begin{example}
We have seen that for $K = \mathbb{Q}(\sqrt{6})$ and $\mathcal{O}_K = \mathbb{Z}[\sqrt{6}]$, $C_K \cong \mathbb{Z}_2$. We will use it to find solutions of the equation:
\[
Y^2 + 54 = X^3 \quad \text{(*)}
\]
\end{example}
\begin{proof} {\bf Claim 1: $2, 3 \nmid y$.}

Suppose on contrary $2 \mid Y$, then obviously $2 \mid X$ and hence $8 \mid X^3 = Y^2 + 54$. But $Y^2+54 \equiv 2 \mod 4$. 

Now suppose $3 \mid Y$. Then $3 | X$. Write $Y = 3a$ and $X = 3b$, one has
$$9a^2 + 54 = 27b^3 \quad \Rightarrow \quad a^2 + 6 = 3b^3$$
which forces $3 \mid a \Rightarrow a =3c$ for some integer $c$. Thus
$$(3c)^2 + 6 = 3b^3 \quad \Rightarrow\quad 3c^2 + 2 = b^3.$$
But $b^3 \equiv 0,1,8 \mod 9$, while $3c^2 + 2 \equiv 2, 5 \mod 9$.  \qed

\medskip
\noindent {\bf Claim 2:
$(Y + 3\sqrt{6}), (Y - 3\sqrt{6})$ has no common prime factors.}

Suppose on contrary $P \mid (Y \pm 3\sqrt{6})$, then as before
$$2 \cdot 3\sqrt{6} \in P \ \Rightarrow\ P \mid (6\sqrt{6}) = (\sqrt{6})^3 \ \Rightarrow \ P \mid (\sqrt{6})\ \Rightarrow \ (\sqrt{6}) \subseteq P \ \Rightarrow \ \sqrt{6} \in P$$
Thus
$$\sqrt{6}, Y + 3\sqrt{6} \in P\ \Rightarrow\ Y \in P \ \Rightarrow \  P \mid (Y)\ \Rightarrow\ N(P) \mid N((Y)) = Y^2$$
But $P \mid  (\sqrt{6})$ implies that $N(P) \mid N((\sqrt{6})) = 6$ and hence $2,3 \mid Y^2$, contradicting Claim 1. \qed


\bigskip
As in Example 1, one can use Claim 2 and the fact that $|C_K| =2$ to conclude that 
$$(Y + 3\sqrt{6}) = (a + b \sqrt{-6})^3 = (\{a^3-18ab^2\} + \{3a^2b-6b^3\}\sqrt{-6})$$
for some $a,b \in \mathbb{Z}$, which yields
$$Y = \pm a(a^2-18b^2) \quad \quad 3 = \pm 3b(a^2-2b^2)$$
So $b = \pm 1$ and $a = 1$ from the second equation, forcing $Y = \pm 17$ in the first equation. Putting it back to the original equation, one has $X = 7$.

\medskip
In conclusion, the only solutions of the Diophantine equation are
\[
(X,Y) = (7, \pm 17).
\]
\end{proof}
